\documentclass[acmsmall,screen,review,anonymous]{acmart}

\settopmatter{printacmref=false}
\setcopyright{none}
\renewcommand\footnotetextcopyrightpermission[1]{}

\usepackage{booktabs}
\usepackage{flafter}
\usepackage{multirow}
\usepackage{placeins}
\usepackage{xspace}

\newcommand{\tool}{SecAware\xspace}
\newcommand{\prompttsg}{Prompt TSG\xspace}
\newcommand{\todo}[1]{\textbf{[TODO: #1]}}

\acmConference[FSE 2027]{The ACM International Conference on the Foundations
of Software Engineering}{July 12--16, 2027}{Shenzhen, China}
\acmYear{2027}
\copyrightyear{2027}

\begin{document}

\title{Discovering and Confirming Prompt-Side Security Mechanisms in LLM Code
Generation}

\author{Anonymous Author(s)}
\affiliation{%
  \institution{Anonymous Institution}
  \country{}
}

\begin{abstract}
Prompts that appear functionally adequate can still omit or weaken constraints
that affect the security of LLM-generated code. Existing explanation and
perturbation methods can identify suspicious text or associations, but they do
not by themselves establish whether a prompt-side factor can be changed while
preserving the task or whether the assigned change improves security on
held-out tasks. We present \tool, a TSG-constrained causal discovery and
randomized counterfactual confirmation framework. \tool extracts a typed
prompt-side task-security graph, constructs local causal-variable tables for
each security scope, and derives background knowledge without treating graph
structure as causation. It runs observational FCI with a G-square conditional
independence test, evaluates path stability with task-cluster bootstrap
resampling, and freezes selected hypotheses before confirmation. On held-out
tasks, four-arm randomized block experiments estimate assigned-arm,
task-clustered intent-to-treat effects on secure-and-functional and
CWE-specific outcomes. A secondary JCI analysis represents randomized arms as
context variables, and optional RFCI provides a sensitivity backend. The
design separates prompt-side mechanisms from generated-code outcome
provenance and retains failed, null, and conflicting evidence. This
design-stage draft specifies the method and evaluation protocol; numerical
results remain intentionally blank until one frozen paper run is available.
\end{abstract}

\begin{CCSXML}
<ccs2012>
 <concept>
  <concept_id>10011007.10011074.10011099</concept_id>
  <concept_desc>Software and its engineering~Software safety</concept_desc>
  <concept_significance>500</concept_significance>
 </concept>
 <concept>
  <concept_id>10002978.10003029.10011703</concept_id>
  <concept_desc>Security and privacy~Software security engineering</concept_desc>
  <concept_significance>500</concept_significance>
 </concept>
</ccs2012>
\end{CCSXML}

\ccsdesc[500]{Software and its engineering~Software safety}
\ccsdesc[500]{Security and privacy~Software security engineering}
\keywords{large language models, code generation, software security, causal
discovery, randomized experiments, prompt analysis}

\maketitle

\section{Introduction}

LLM-based code generation makes natural-language prompts part of the software
security boundary. A prompt may state the required functionality while leaving
trust boundaries, validation, authorization, or safe-API constraints implicit.
The resulting code can be functional yet insecure. A useful explanation must
therefore do more than rank tokens: it should identify an editable prompt-side
factor, state the uncertainty of its observational evidence, and test a
prespecified intervention without changing the underlying programming task.

Three issues make this difficult. First, prompt requirements have typed and
relational structure that surface text alone may obscure. Second, structure in
a prompt representation must not be mistaken for a causal graph. Third,
confirmation is vulnerable to post-treatment selection if only successful or
semantically convenient edits are analyzed. A defensible workflow must keep
feature extraction, causal discovery, intervention assignment, generated-code
evaluation, and effect estimation separate and auditable.

\tool addresses these issues with a Prompt-only design. It represents the
input as a \prompttsg, translates graph-derived features into local statistical
tables, and constrains rather than replaces a standard FCI implementation.
Candidate paths are stability-screened and frozen before a randomized,
held-out confirmation stage. Generated code never supplies causal mechanism
features: it is consumed only by an independent security Oracle and functional
evaluator. Confirmation uses the assigned arm as treatment and includes all
assigned units under a prespecified failure policy.

This paper makes the following contributions:

\begin{itemize}
  \item a Prompt-only typed representation and provenance-separated conversion
  to local causal-variable tables for individual CWE/security scopes;
  \item a constrained observational discovery protocol using causal-learn FCI,
  G-square tests, TSG-derived background knowledge, and task-clustered
  stability analysis;
  \item a prevalidated four-arm randomized block design with task-clustered ITT
  estimation, diagnostic treatment-fidelity measures, and secondary JCI and
  RFCI analyses; and
  \item an evaluation plan that compares discovery-to-confirmation yield,
  representation and analysis ablations, defensive interventions, and expert
  ratings without converting exploratory evidence into quantitative claims.
\end{itemize}

\todo{Replace this note with a security-neutral motivating prompt, its frozen
hypothesis, four assigned prompt variants, and outcome summary from the final
paper run.}

\section{Problem Setup and Scope}

\subsection{Units, Splits, and Outcomes}

Let a generation unit be indexed by task $t$, prompt variant $v$, model $m$,
and randomized seed slot $s$:

\[
  u = (t, v, m, s).
\]

Tasks are split before model-based discovery into disjoint discovery and
confirmation pools. The discovery pool supports observational hypothesis
formation. The held-out confirmation pool supports prompt construction,
randomized assignment, code generation, and effect estimation. Task identity
is the clustering unit because variants and seed slots from the same task are
not independent.

The primary outcome is \emph{secure-and-functional}: generated code must pass
the locked functional contract and the security Oracle. A CWE-specific security
outcome is secondary. The run manifest defines fail-closed handling for
generation, execution, and Oracle failures before assignment; these outcomes
are not chosen after observing an arm.

\subsection{Causal Variables}

For each CWE/security scope and model, \tool constructs a separate local table
with four variable groups:

\begin{itemize}
  \item $W$: pre-treatment task and generation metadata allowed by the locked
  analysis specification;
  \item $X$: direct Prompt TSG feature states and relational Prompt TSG motifs;
  \item $Y$: independently produced secure-and-functional and CWE-specific
  outcomes; and
  \item $C$: randomized arm-context variables used only in the secondary JCI
  analysis.
\end{itemize}

The Prompt TSG and the learned partial ancestral graph (PAG) have different
semantics. Prompt TSG edges encode typed task, security, and presentation
relations. PAG endpoint marks summarize independence information under FCI.
An edge in the Prompt TSG is never reported as a causal edge. The generated
program is likewise not a causal-variable graph: it is an evaluation object for
the Oracle and functional evaluator.

\subsection{Claim Boundary}

The framework does not invent a new general-purpose causal discovery
algorithm, identify a unique causal DAG, or interpret every possible PAG path
as a mechanism. Latent confounding is allowed by FCI, and observational output
is labeled as candidate evidence. Paper-facing confirmation claims require a
frozen hypothesis and a randomized assigned-arm contrast on held-out tasks.

\section{Methodology}

\begin{figure}[t]
  \centering
  \fbox{\parbox{0.92\linewidth}{\centering
    \vspace{1.2em}
    Discovery tasks $\rightarrow$ Prompt TSGs $\rightarrow$ local tables and
    background knowledge $\rightarrow$ FCI/bootstrap $\rightarrow$ frozen
    hypotheses $\rightarrow$ prevalidated four-arm prompts $\rightarrow$
    randomized held-out generation $\rightarrow$ Oracle/functional outcomes
    $\rightarrow$ task-clustered ITT and JCI sensitivity
    \vspace{1.2em}}}
  \caption{The \tool discovery and randomized confirmation workflow. The final
  paper will replace this placeholder with a provenance-aware pipeline figure.}
  \Description{A pipeline from discovery tasks through Prompt TSG extraction,
  constrained FCI, frozen hypotheses, randomized four-arm confirmation,
  generated-code evaluation, task-clustered ITT estimation, and JCI analysis.}
  \label{fig:overview}
\end{figure}

\subsection{Prompt TSG Extraction and Construction}

The Prompt TSG unifies task, security, and presentation features in one typed
graph. Task nodes encode requested operations and functional constraints;
security nodes encode trust assumptions, assets, guards, sanitization
requirements, and security-relevant relations; presentation nodes encode
wording or formatting properties used for placebo control. This unified view
allows an intervention to name its target and its permitted delta while keeping
non-target layers invariant.

The implementation supports three run-locked extraction/construction
strategies:

\begin{enumerate}
  \item an LLM emits schema-constrained semantic facts and a deterministic
  builder creates the graph;
  \item an LLM emits a complete typed graph that is schema-validated and
  canonicalized; and
  \item a finite catalog/rule extractor emits facts for deterministic graph
  construction as a reproducible fallback.
\end{enumerate}

The paper run freezes the selected strategy, model, system prompt, schema,
catalog, and canonicalization version. Extractors do not receive generated
code, Oracle findings, or outcome labels. Invalid extraction is recorded as an
explicit failure rather than repaired with an expanding collection of hidden
security rules.

\subsection{Local Tables and TSG-Derived Background Knowledge}

The builder projects only variables relevant to one CWE/security scope and one
model into each analysis table. It records variable type, state domain,
missingness provenance, task cluster, and source artifact. Relational motifs
are finite, versioned functions of the Prompt TSG, such as a security
requirement being scoped to a sensitive operation. They describe prompt
structure, not code-side data flow.

The same typed schema derives background knowledge for discovery:

\begin{itemize}
  \item temporal tiers prevent outcomes from causing pre-treatment metadata or
  prompt features;
  \item forbidden directions encode provenance and time-order constraints; and
  \item typed adjacency restrictions remove semantically impossible local
  adjacencies.
\end{itemize}

These constraints do not require the $X$--$Z$ or $Z$--$Y$ edges of a candidate
path. Requiring a candidate path in background knowledge would turn discovery
into confirmation of a supplied answer.

\subsection{Observational FCI and Stability}

The minimum runnable backend is causal-learn FCI with the G-square conditional
independence test. It requires no Java runtime. For the reference observational
draw, the pipeline randomly selects one seed per task so that repeated
generations do not masquerade as independent observations. It then reruns the
analysis across task-cluster bootstrap samples, resampling entire tasks and
selecting one seed per sampled task according to the frozen policy.

Each run preserves the raw PAG, background-knowledge digest, CI-test
configuration, selected rows, random seed, and backend diagnostics. The path
extractor searches for endpoint-compatible possible $X$--$Z$--$Y$ paths and
summarizes adjacency and endpoint stability across bootstrap replicates. It
does not orient circles merely to make a path directional. A candidate is
eligible for freezing only if it meets prespecified support, stability,
targetability, and provenance conditions.

Before confirmation begins, \tool freezes the selected top-$K$ hypotheses,
their native evidence, feature/TargetSpec mappings, expected directions,
security scopes, models, operation families, and analysis code/configuration.
No confirmation outcome may change candidate rank or mapping.

\subsection{Intervention Construction and Prevalidation}

An intervention can be executed text-natively or graph-natively. A text-native
executor edits the original prompt under a TargetSpec. A graph-native executor
applies a typed graph patch and verbalizes the resulting Prompt TSG. Either path
may use a deterministic renderer or a locked LLM executor. The final run chooses
one configuration in advance and records all prompts, model parameters,
schemas, and retries. Treating the LLM executor as an ideal implementation
assumption does not remove the need to measure realized fidelity as a
diagnostic.

All variants are constructed, validated, and frozen before randomized code
generation. Hard pre-randomization gates check schema validity, provenance,
security-neutral task preservation, family-specific AllowedDelta constraints,
and absence of outcome or vulnerability-label leakage. A safety intervention
may change only its target security feature; a functionality intervention may
change only its target task feature; a placebo may change only presentation.
Invalid blocks are rejected before assignment and retained in failure
artifacts.

After assignment, re-extraction records whether the target changed, whether
the prompt remained semantically compliant, and whether non-target features
drifted. These variables diagnose intervention execution. They are not primary
ITT eligibility filters.

\subsection{Randomized Confirmation Arms}

For a safety-ADD hypothesis, each eligible block contains four arms:

\begin{itemize}
  \item \verb|TARGET_PATCH|: add the frozen target security feature;
  \item \verb|NOOP_REWRITE|: rewrite without changing task or security state;
  \item \verb|LENGTH_MATCHED_PLACEBO|: make a presentation-only edit matched
  to the target edit's length; and
  \item \verb|GENERIC_SECURITY_REMINDER|: add a non-targeted security reminder.
\end{itemize}

For safety-REMOVE, the analogous arms are:

\begin{itemize}
  \item \verb|TARGET_REMOVE|;
  \item \verb|NOOP_RETAIN|;
  \item \verb|LENGTH_MATCHED_SHAM_EDIT|; and
  \item \verb|GENERIC_SECURITY_REPLACEMENT|.
\end{itemize}

ADD and REMOVE are distinct protocols and are never pooled into one effect.
Within each locked task/hypothesis/TargetSpec/protocol/model block, randomized
seed slots are balanced across arms. Assignment artifacts are frozen before
generation.

The target-versus-no-op contrast estimates the effect of the specified feature
change relative to an operation-matched rewrite. The placebo contrast checks
whether wording/length alone explains the effect. The generic-reminder contrast
separates target-specific content from a broad security cue.

\subsection{Outcome Construction and Provenance Separation}

Each assigned prompt is sent to the locked generation model. Generated code is
evaluated by an independent, versioned security Oracle and functional
evaluator. The run manifest freezes tool versions, policies, timeouts, and
failure handling. Outcome artifacts may reference prompt, assignment, code,
Oracle, and functional-evaluation digests, but the prompt extractor and
hypothesis freezer cannot read downstream labels. This one-way provenance
boundary prevents outcome leakage into mechanism construction.

\subsection{Primary ITT Estimation}

For frozen hypothesis $h$, target $q$, operation $o$, model $m$, and outcome
$Y$, the primary estimand is the assigned-arm risk difference

\[
  \Delta^{\mathrm{ITT}}_{h,q,o,m}
  = \mathbb{E}[Y \mid A=\mathrm{target}]
  - \mathbb{E}[Y \mid A=\mathrm{no\mbox{-}op}].
\]

The assigned arm $A$ is used regardless of realized target change or semantic
diagnostics. The denominator therefore includes every randomized unit under
the prespecified outcome/failure policy. Uncertainty is clustered by task, and
the multiplicity procedure is frozen for each registered hypothesis family.
Effects are not pooled across hypothesis, target, operation, model, or CWE
unless a separate pooled estimand is preregistered.

Target-versus-placebo and target-versus-generic contrasts are secondary
specificity analyses. Per-protocol summaries may be reported only as clearly
labeled diagnostics. Task-level improvement, harm, and tie rates can describe
outcome transitions, but they are not interpreted as individual causal flips.

\subsection{JCI and RFCI Sensitivity Analyses}

In the secondary Joint Causal Inference (JCI) analysis, arm indicators enter as
context variables $C$. \tool first stores the raw augmented PAG, then applies
only the frozen JCI assumptions and stores the constrained PAG plus every
orientation delta. Deterministic relations among one-hot arm indicators are
handled by the locked JCI encoding and synthetic tests; they are not silently
treated as ordinary independent features.

An optional py-tetrad RFCI backend tests sensitivity to the discovery backend.
It is not required for the minimum Python-only pipeline. Any violation of
background-knowledge or endpoint contracts is surfaced as a backend failure,
not hidden by post-hoc orientation.

\subsection{Evidence and Failure Artifacts}

The reporting layer preserves PAG, hypothesis, intervention, assignment,
effect, diagnostic, outcome, and failure artifacts with configuration and
content digests. It uses the following evidence vocabulary:

\begin{enumerate}
  \item Observational Candidate;
  \item Confirmed Intervention Effect;
  \item Target-Specific Effect;
  \item Bidirectional Support;
  \item Directionally Consistent but Inconclusive; and
  \item Null / Conflicting / Non-Evaluable.
\end{enumerate}

The distinction keeps observational paths, randomized effects, specificity,
and reverse-operation support from collapsing into one binary label.

\section{Research Questions and Evaluation Design}

\subsection{RQ1: Discovery and Confirmation}

\textbf{RQ1. How effectively can different methods discover and confirm
security-relevant prompt-side mechanisms in LLM code generation?}

Methods receive the same discovery pool, held-out confirmation pool, fixed
native candidate budget $K$, model/seed policy, Oracle, randomization protocol,
and ITT estimator. Each method first emits native candidates without seeing
confirmation outcomes. Candidates are then mapped to the finite
FeatureSpec/TargetSpec catalog and frozen before prompt construction.

The comparison follows a fixed funnel:

\[
K \rightarrow N_{\mathrm{native}} \rightarrow N_{\mathrm{mapped}}
\rightarrow N_{\mathrm{protocol}} \rightarrow N_{\mathrm{randomized}}
\rightarrow N_{\mathrm{confirmed}}.
\]

Candidate utilization, mapping coverage, mapped yield, protocolization rate,
block-freeze coverage, and randomized yield reveal where a method loses usable
hypotheses. The method-level primary metric is confirmed yield at $K$.
Conditional confirmation among randomized hypotheses is diagnostic because it
omits earlier funnel losses. Hypothesis-specific task-clustered ITT estimates
retain direction, interval, and multiplicity-adjusted status.

\begin{table}[htbp]
  \centering
  \small
  \caption{RQ1 fixed-$K$ discovery-to-confirmation funnel. All cells remain
  blank until the frozen paper run.}
  \label{tab:rq1}
  \begin{tabular}{lrrrrrr}
    \toprule
    Method & $K$ & Native & Mapped & Protocol & Rand. & Confirmed \\
    \midrule
    \tool & -- & -- & -- & -- & -- & -- \\
    Native baseline 1 & -- & -- & -- & -- & -- & -- \\
    Native baseline 2 & -- & -- & -- & -- & -- & -- \\
    Native baseline 3 & -- & -- & -- & -- & -- & -- \\
    \bottomrule
  \end{tabular}
\end{table}
\FloatBarrier

\subsection{RQ2: Representation and Causal Analysis}

\textbf{RQ2. How do SecAware's structured representation and causal analysis
components contribute to mechanism discovery and confirmation?}

RQ2 uses a preregistered two-by-two ablation. The full and association variants
share direct and relational Prompt TSG motifs. The reduced variants retain only
direct catalog feature states. FCI variants use the complete constrained
discovery protocol; association variants rank candidates with univariate
G-square association under the same candidate budget. All downstream mapping,
freezing, randomization, and ITT steps remain identical.

\begin{table}[htbp]
  \centering
  \small
  \caption{RQ2 two-by-two component design and placeholder outcomes.}
  \label{tab:rq2}
  \begin{tabular}{llllr}
    \toprule
    Variant & Representation & Selection & Rand. & Confirmed \\
    \midrule
    Full & direct + relational & FCI & -- & -- \\
    Reduced repr. & direct only & FCI & -- & -- \\
    Assoc. selection & direct + relational & univ. G-square & -- & -- \\
    Double ablation & direct only & univ. G-square & -- & -- \\
    \bottomrule
  \end{tabular}
\end{table}
\FloatBarrier

The primary comparison uses confirmed yield at the same $K$; funnel losses and
hypothesis-specific ITT effects explain whether a difference arises from
representation, discovery, operationalization, or confirmation.

\subsection{RQ3: Defensive Interventions}

\textbf{RQ3. Which prompt-side defensive interventions effectively reduce
insecure code generation?}

RQ3 evaluates frozen defensive targets separately under safety-ADD and
safety-REMOVE protocols. ADD asks whether introducing a specific defensive
requirement improves secure-and-functional outcomes. REMOVE asks whether
removing that requirement worsens the outcome and can provide bidirectional
support. Neither protocol selects tasks using observed baseline insecurity.

\begin{table}[htbp]
  \centering
  \small
  \caption{RQ3 intervention-effect skeleton. Primary contrasts are target
  versus operation-matched no-op assigned-arm ITT effects.}
  \label{tab:rq3}
  \begin{tabular}{lllrrl}
    \toprule
    Target & Op. & Outcome & ITT RD & Interval & Evidence \\
    \midrule
    Input validation & ADD & secure+func. & -- & -- & -- \\
    Input validation & REMOVE & secure+func. & -- & -- & -- \\
    Safe API constraint & ADD & secure+func. & -- & -- & -- \\
    Safe API constraint & REMOVE & secure+func. & -- & -- & -- \\
    Authorization guard & ADD & secure+func. & -- & -- & -- \\
    Authorization guard & REMOVE & secure+func. & -- & -- & -- \\
    \bottomrule
  \end{tabular}
\end{table}
\FloatBarrier

Target-versus-placebo and target-versus-generic contrasts assess specificity.
Target-change, semantic-compliance, non-target-drift, task improvement, harm,
and tie rates are reported as execution and transition diagnostics rather than
filters or individual-effect claims.

\subsection{RQ4: Expert-Perceived Explanation Utility}

\textbf{RQ4. How do security experts rate and rank the perceived quality and
usefulness of explanations produced by different methods?}

RQ4 is a subjective, conditional-utility study. Each method's native output is
placed in a neutral wrapper with the same fields: Finding, Evidence, Expected
direction, Experimental support or uncertainty, and Suggested audit or repair
action. Method names are hidden and case/method order is balanced.

Experts rate clarity, evidence sufficiency, credibility, actionability, and
overall usefulness. Overall usefulness is the preregistered primary dimension;
the remaining dimensions explain the ranking. The analysis uses an ordinal
mixed-effects model with method as a fixed effect and participant and case as
random effects, followed by multiplicity-adjusted comparisons. The RQ does not
claim objective mechanism-identification or repair accuracy.

\begin{table}[htbp]
  \centering
  \small
  \caption{RQ4 expert-rating skeleton. Entries will report ordinal summaries
  and model-based contrasts from the completed study.}
  \label{tab:rq4}
  \begin{tabular}{lrrrrr}
    \toprule
    Method & Clarity & Evidence & Credibility & Actionability & Usefulness \\
    \midrule
    \tool & -- & -- & -- & -- & -- \\
    Baseline 1 & -- & -- & -- & -- & -- \\
    Baseline 2 & -- & -- & -- & -- & -- \\
    Baseline 3 & -- & -- & -- & -- & -- \\
    \bottomrule
  \end{tabular}
\end{table}
\FloatBarrier

\section{Synthetic Validation}

Synthetic structural causal models test the analysis contracts independently
of the paper dataset. The suite includes a true $X$--$Z$--$Y$ chain, latent
confounding, a null factor, and deterministic JCI context relations. Tests check
PAG/path compatibility, bootstrap stability behavior, preservation of raw and
constrained JCI outputs, and fail-visible backend behavior. Synthetic recovery
validates implementation logic; it is not evidence that a paper hypothesis is
true.

\section{Implementation and Study Status}

The repository implements the core Prompt-only pipeline: Prompt TSG extraction
and canonicalization, local tables and background knowledge, causal-learn FCI,
task-cluster bootstrap stability, hypothesis freezing, safety intervention
families, randomized blocks, task-clustered ITT estimation, JCI artifacts,
optional RFCI sensitivity, and the independent Oracle/functional-evaluation
path.

This status does not imply that all paper RQs are complete. External
method-native adapters and mapping audits, offline RQ2 ablation runs, final RQ1
and RQ3 table builders and summaries, the RQ4 study infrastructure and expert
study, and the final frozen paper configuration remain to be completed. Current
demo/mock \texttt{paper\_v0} artifacts are development evidence and must not be
used as final paper results. Consequently, this draft contains no quantitative
findings.

\section{Threats to Validity}

\paragraph{Construct validity.}
Security and functionality outcomes depend on the locked Oracle and test
coverage. Static or dynamic tools can miss vulnerabilities, misclassify safe
idioms, or fail on unsupported code. The final study must publish tool versions,
policies, calibration audits, disagreement handling, and failure denominators.

\paragraph{Discovery validity.}
FCI returns equivalence-class information rather than a unique causal graph.
Finite samples, sparse discrete tables, latent confounding, and background
knowledge can alter PAG endpoints. We preserve circle endpoints, raw PAGs,
bootstrap instability, and null/conflicting candidates instead of resolving
ambiguity narratively.

\paragraph{Internal validity.}
LLM prompt executors can introduce noncompliance or non-target drift. Freezing
all variants before generation and randomizing assigned arms protects the
primary contrast, while treatment-fidelity diagnostics reveal dilution.
Analyzing only successful edits would reintroduce post-treatment selection, so
the primary ITT denominator is assignment-based.

\paragraph{Multiplicity and selection.}
Candidate ranking, top-$K$ selection, multiple hypotheses, outcomes, models,
and contrasts create multiplicity. The final run freezes families, directions,
and adjustments before confirmation and keeps exploratory analyses separate.

\paragraph{External validity.}
Effects can vary across model families, languages, CWE scopes, task sources,
prompt styles, and generation settings. Results are reported within the locked
scope rather than generalized to all LLM code generation.

\paragraph{Human-study validity.}
RQ4 measures expert perceptions conditional on selected cases and a shared
wrapper. Ratings may depend on participant experience, explanation length, and
interface order. Blinding, balanced presentation, participant/case random
effects, and transparent exclusions limit but do not eliminate these threats.

\section{Reproducibility Plan}

The anonymous replication package will contain frozen discovery and
confirmation manifests; task splits; Prompt TSG and local-table artifacts;
background-knowledge digests; raw and constrained PAGs; bootstrap records;
frozen hypotheses and mappings; prompt variants and validation results;
assignment blocks; generated-code digests; Oracle and functional outcomes;
task-clustered effect outputs; JCI/RFCI diagnostics; failure artifacts; and
scripts that rebuild every results table. Each artifact will record schema,
configuration, code, and parent-content digests. Exploratory and demo runs will
remain outside the paper table pipeline.

\section{Conclusion}

\tool combines Prompt TSG-constrained FCI with randomized held-out confirmation
to distinguish observational prompt-side candidates from assigned-arm security
effects. Its design keeps prompt structure separate from causal edges, code
evaluation separate from mechanism extraction, and intervention diagnostics
separate from the primary task-clustered ITT denominator. The approved
evaluation asks how methods traverse the discovery-to-confirmation funnel, how
representation and causal analysis contribute, which defensive interventions
work, and how experts perceive the resulting explanations. Quantitative claims
await one frozen paper run.

\section{Data Availability}

This design-stage manuscript reports no quantitative results and therefore has
no frozen result dataset to release. Before a results-bearing submission, we
will archive the anonymous replication package described above and replace this
statement with its access conditions and persistent location. No demo or mock
artifact will be presented as final paper evidence.

\end{document}
